\chapter{Alıştırmaların Kısa Çözümleri}

Bu ekte, bölüm sonlarındaki \emph{Alıştırmalar} için kısa çözüm yolları verilmiştir. Birden fazla doğru yaklaşımın mümkün olduğu açık uçlu sorularda örnek bir çözüm veya değerlendirme ölçütü sunulmuştur. Alıştırma kısmı içeren bölümler aşağıda kitaptaki sırasıyla yer alır; diğer bölümlerin pekiştirme çalışması gözden geçirme soruları ve çoktan seçmeli sorular üzerinden yürütülür.

\section{Bölüm~\ref{chap:muhakeme-turleri}: Muhakeme Türleri}

\begin{enumerate}
\item Bir gazete başlığında çoğu kez sınırlı gözlemlerden genel sonuca gidildiği için tümevarım; matematiksel bir iddianın aksiyom ve tanımlardan çıkarılmasında tümdengelim; gözlenen veriyi en iyi açıklayan bilimsel hipotezin kurulmasında abdüksiyon kullanılır. Gerçek örneklerde birden fazla muhakeme türü birlikte bulunabilir.

\item $n=1,2,3,4,5$ için değerler sırasıyla $43,47,53,61,71$ olup asaldır. Bu hesaplar iddiayı destekler, fakat bütün doğal sayılar için ispatlamaz. Nitekim $n=40$ için $n^2+n+41=41^2$, $n=41$ için de $n^2+n+41=41\cdot43$ olur; iki değer de bileşiktir.

\item Geçerli fakat sağlam olmayan örnek: ``Bütün kareler çemberdir; bütün çemberler çokgendir; öyleyse bütün kareler çokgendir.'' Çıkarım biçimi geçerli, öncüller yanlıştır. Sağlam örnek: ``Dördün katı olan her tam sayı çifttir; 12 dördün katıdır; öyleyse 12 çifttir.''

\item (a) Geçersizdir; sonucu onaylama hatasıdır. (b) Geçerlidir; modus tollens kalıbıdır. (c) Geçersizdir; öncülü inkâr hatasıdır. Eşkenar olmayan bir üçgen yine de ikizkenar olabilir.

\item Wason görevinde doğru seçim, ünlü harf bulunan kart ile tek sayı bulunan karttır: kuralı bozabilecek durumlar bunların arka yüzlerindedir. Uygulama yanıtları değişebilir; değerlendirmede sonucu onaylama ve karşı örneği aramama eğilimleri özellikle kaydedilmelidir.

\item Genel konumda en çok bölge sayısı
\[
1+\binom{n}{2}+\binom{n}{4}
\]
olur. $n=6$ için $1+15+15=31$ elde edilir. İlk beş değer $1,2,4,8,16$ olduğundan örüntü ikinin kuvvetleri izlenimi verir; altıncı adım bu yüzeysel genellemeyi bozar.

\item Pisagor teoreminde dik üçgenin hipotenüsü üzerindeki karenin alan karesi, dik kenarlar üzerindeki karelerin alan kareleri toplamına karşılık gelir. De Gua teoreminde dik köşeli dörtyüzlünün karşı yüz alanının karesi, dik köşede birleşen üç yüzün alan kareleri toplamıdır. Analoji kenarı yüze, üçgeni dörtyüzlüye ve iki dik bileşeni üç dik bileşene eşler; bu eşleştirme benzerliğin yüzeysel değil yapısal olduğunu gösterir.

\item Naif empirik gerekçe birkaç çarpımı denemektir. Genel örnekte $n(n+1)(n+2)$ yazılır ve üç ardışık sayıdan birinin çift, birinin 3'ün katı olduğu görülür. Tümdengelimsel gerekçede bu iki bölünebilme olgusu sistemli biçimde kullanılarak çarpımın hem 2'ye hem 3'e, dolayısıyla 6'ya bölündüğü kanıtlanır.
\end{enumerate}

\section{Bölüm~\ref{chap:muhakeme-ogretimi}: Düşünme Düzeyleri ve Muhakeme Öğretimi}

\begin{enumerate}
\item Örnek olarak 2. düzeydeki öğrencilere farklı yönlerde çizilmiş dörtgen kartları verilebilir. Öğrenciler şekilleri kenar, açı, paralellik ve köşegen özelliklerine göre sınıflandırır; görünüş yerine özellik kullanmaları istenir. Etkinlik, çözümleme düzeyinin karakteristik işi olan özellikleri belirleme ve karşılaştırmaya uygundur.

\item Öğrenciye ``Tek bir örnek seçebilir misin; bu örnekte kullandığın gerekçenin sayı değişse de neden değişmediğini gösterebilir misin?'' sorusu yöneltilebilir. Böylece hesaplanan özel sayıdan, bütün durumları temsil eden $2k+1$ gibi genel bir örneğe geçmesi beklenir.

\item Yanıt seçilen ispata göre değişir. Tam puan için ispatta sonradan görünmez hâle gelen en az bir keşif kararı (yardımcı çizgi, temsil, lemma veya yöntem seçimi) belirtilmeli ve öğrenciyi bu karara götürecek açık bir öğretmen sorusu yazılmalıdır.

\item ``Kare dikdörtgen değildir'' sözü genellikle özellikleri listeleyebilen fakat sınıf içerme ilişkisini kuramayan 2. van Hiele düzeyine işaret eder. Örnek soru dizisi: ``Bir şeklin dikdörtgen olması için hangi özellikler yeterlidir?'', ``Kare bu özelliklerin her birini sağlıyor mu?'', ``Bir şeklin ek özellik taşıması onun dikdörtgen olmasını ortadan kaldırır mı?'' Bu sorular öğrenciyi 3. düzeydeki kapsayıcı sınıflamaya taşır.
\end{enumerate}

\section{Bölüm~\ref{chap:tumevarim}: Tümevarımla İspat}

\begin{enumerate}
\item $n=1$ için iki taraf da $2$'dir. $\sum_{j=1}^{k}j(j+1)=k(k+1)(k+2)/3$ varsayılsın. Buna $(k+1)(k+2)$ eklenirse
\[
\frac{k(k+1)(k+2)}{3}+(k+1)(k+2)
=\frac{(k+1)(k+2)(k+3)}{3}
\]
elde edilir.

\item Taban adımı $n=5$ için $32>25$'tir. $2^k>k^2$ varsayılırsa $2^{k+1}>2k^2$. $k\geq5$ için $2k^2>(k+1)^2$ olduğundan sonuç $k+1$ için de geçerlidir.

\item Genel konumdaki $n$'inci doğru önceki doğrularla $n-1$ farklı noktada kesilir ve $n$ parçaya ayrılarak $n$ yeni bölge oluşturur. Bu nedenle $R_n=R_{n-1}+n$, $R_0=1$; toplamdan $R_n=1+n(n+1)/2$ bulunur.

\item $F_1=F_2=1$ için iddia açıktır. $F_j<2^j$ eşitsizliğinin $j\leq k$ için doğru olduğunu varsayalım. $F_{k+1}=F_k+F_{k-1}<2^k+2^{k-1}<2^{k+1}$ olur. İki önceki terim kullanıldığı için güçlü tümevarım doğaldır.

\item Hata $n=1$'den $n=2$'ye geçiştedir. İki atlı grupta ``ilk $k$ at'' ile ``son $k$ at'' gruplarının ortak bir atı yoktur; bu nedenle renk eşitliğini bir gruptan diğerine taşıyan bağlantı kurulamaz.
\end{enumerate}

\section{Bölüm~\ref{chap:kumeler}: Kümeler ve Küme Üzerine Muhakeme}

\begin{enumerate}
\item Her $x$ için gösterge fonksiyonlarını $0$ ve $1$ değerli alalım. $\chi_{A\triangle B}(x)=\chi_A(x)+\chi_B(x)$ eşitliği mod 2'de geçerlidir. Toplamanın birleşme özelliğinden $\chi_{(A\triangle B)\triangle C}=\chi_{A\triangle(B\triangle C)}$ elde edilir; dolayısıyla simetrik fark birleşmelidir.

\item $A\times B$ kümesinde $mn$ sıralı ikili vardır. Her ikilinin bağıntıda bulunup bulunmaması bağımsız iki seçenek verdiğinden bağıntı sayısı $2^{mn}$'dir. $A$'dan $B$'ye her fonksiyon, $A$'nın her elemanı için $n$ seçenek verdiğinden fonksiyon sayısı $n^m$'dir.

\item $X\in\mathcal P(A)\cup\mathcal P(B)$ ise $X\subseteq A$ veya $X\subseteq B$ olur; her iki durumda da $X\subseteq A\cup B$ olduğundan $X\in\mathcal P(A\cup B)$. Eşitlik genelde yanlıştır: $A=\{1\}$ ve $B=\{2\}$ için $\{1,2\}\in\mathcal P(A\cup B)$, fakat sol tarafta değildir.

\item Kesişim $\{0\}$'dır. Sıfır bütün aralıklardadır. $x>0$ için $n>1/x$ seçilirse $1/n<x$ olur; dolayısıyla $x$, $[0,1/n]$ aralığında bulunmaz.

\item ``Kendisini tıraş etmeyen bütün ve yalnız erkekleri tıraş eden berber'' kendisini tıraş ederse tanıma göre etmemeli, etmezse tanıma göre etmelidir. Russell kümesindeki $R\in R\Longleftrightarrow R\notin R$ çelişkisiyle aynı öz-gönderimli yapı görülür; fark, berber örneğinde böyle bir berberin bulunmadığı söylenebilirken naif küme kuramında kapsam ilkesinin sınırlandırılması gerekmesidir.
\end{enumerate}

\section{Bölüm~\ref{chap:bagintilar}: Bağıntılar ve Denklik}

\begin{enumerate}
\item $(x_1,y_1)\sim(x_2,y_2)$ bağıntısını $x_1^2+y_1^2=x_2^2+y_2^2$ ile tanımlayalım. Eşitliğin yansıma, simetri ve geçiş özellikleri bağıntının denklik olduğunu gösterir. $r\geq0$ için denklik sınıfı, merkezî orijinde ve yarıçapı $r$ olan çemberdir; $r=0$ sınıfı yalnızca orijinden oluşur.

\item Bölünebilme sıralamasının Hasse örtmeleri şunlardır: $1$--$2$, $1$--$3$, $2$--$4$, $2$--$6$, $3$--$6$, $4$--$12$ ve $6$--$12$. Diyagramda $1$ altta, $12$ üstte yer alır.

\item Eşitlik bağıntısı hem simetrik hem ters simetriktir: $a=b$ ise $b=a$; ayrıca $a=b$ ve $b=a$ koşulları zaten $a=b$ sonucunu verir.

\item Her $a$, yansıma nedeniyle $[a]$ sınıfındadır; dolayısıyla sınıfların birleşimi bütün kümeyi örter. $[a]\cap[b]$ boş değilse bir $x$ için $x\sim a$ ve $x\sim b$ olur. Simetri ve geçişlilikten $a\sim b$; buradan $[a]=[b]$ çıkar. Böylece sınıflar ya eşit ya ayrıktır ve bir bölümleme oluşturur.

\item Mod 5 denklik sınıfları $[0],[1],[2],[3],[4]$'tür; burada $[r]=\{5k+r:k\in\mathbb Z\}$. Her tam sayı bölme algoritmasıyla bu sınıflardan tam birine girer.
\end{enumerate}

\section{Bölüm~\ref{chap:kongruanslar}: Kongrüanslar}

\begin{enumerate}
\item $7\equiv2$ (mod 5) ve $2^4\equiv1$ (mod 5). $100=4\cdot25$ olduğundan $7^{100}\equiv1$ (mod 5).

\item $3\cdot5=15\equiv1$ (mod 7) olduğundan $x\equiv5$ (mod 7).

\item $10\equiv-1$ (mod 11) olduğundan $a_k10^k\equiv a_k(-1)^k$ (mod 11). Bu yüzden bir sayının 11'e göre kalanı, sağdan başlayarak rakamlarının dönüşümlü toplamının kalansını verir; toplam 11'in katıysa sayı da 11'in katıdır.

\item $36=2^2\cdot3^2$ olduğundan
\[
\varphi(36)=36\left(1-\frac12\right)\left(1-\frac13\right)=12.
\]

\item $p=5$, $q=11$ için $n=pq=55$ ve $\varphi(n)=4\cdot10=40$'tır. Örneğin $e=3$ seçilebilir. $3d\equiv1$ (mod 40) koşulunu $d=27$ sağlar. Böylece açık anahtar $(55,3)$, özel üs $27$ olur.
\end{enumerate}

\section{Bölüm~\ref{chap:fonksiyonlar}: Fonksiyonlar}

\begin{enumerate}
\item $f(x)=x/(1+|x|)$ için $x\geq0$ iken $f(x)=x/(1+x)$, $x<0$ iken $f(x)=x/(1-x)$ olur. Her iki parça kesin artandır ve görüntüler birlikte $(-1,1)$'i verir. Ters fonksiyon
\[
f^{-1}(y)=\frac{y}{1-|y|}\qquad(-1<y<1)
\]
biçimindedir; bu da birebirlik ve örtenliği doğrular.

\item $f:A\to B$ birebirse $|f(A)|=|A|$ olur. Sonlu kümelerde $|A|=|B|$ olduğundan $|f(A)|=|B|$ ve $f(A)\subseteq B$ eşitliği ancak $f(A)=B$ iken mümkündür; dolayısıyla $f$ örtendir.

\item $y\in f(A_1\cup A_2)$ ancak ve ancak bir $x\in A_1\cup A_2$ için $y=f(x)$ olur. Bu da $x\in A_1$ veya $x\in A_2$, yani $y\in f(A_1)\cup f(A_2)$ demektir. Böylece iki küme eşittir.

\item $g\circ f$ birebir ve $f(a_1)=f(a_2)$ olsun. Bu durumda $g(f(a_1))=g(f(a_2))$; bileşkenin birebirliğinden $a_1=a_2$ çıkar, dolayısıyla $f$ birebirdir. Buna karşılık $A=\{0,1\}$, $B=\{0,1,2\}$ için $f(0)=0$, $f(1)=1$; $g(0)=0$, $g(1)=1$, $g(2)=1$ alınırsa $g\circ f$ birebir, fakat $g$ birebir değildir.

\item Her $x$ için $\chi_{A\triangle B}(x)=1$ tam olarak $x$'in $A$ ile $B$'den yalnız birinde bulunması durumunda gerçekleşir. Bu, mod 2'de $\chi_A(x)+\chi_B(x)$ değeridir; dolayısıyla $\chi_{A\triangle B}=\chi_A+\chi_B$ (mod 2).
\end{enumerate}

\section{Bölüm~\ref{chap:kardinalite}: Kardinalite ve Sonsuzluk}

\begin{enumerate}
\item $\mathbb N\times\mathbb N$ çiftleri $m+n=0,1,2,\ldots$ köşegenleri boyunca sıralanabilir. Açık bir numaralama örneği
\[
\pi(m,n)=\frac{(m+n)(m+n+1)}{2}+n
\]
fonksiyonudur. Bu fonksiyon bir eşleme verdiğinden çarpım kümesi sayılabilirdir.

\item $|A|=n$ ise her eleman bir alt kümeye alınır veya alınmaz; bu nedenle $|\mathcal P(A)|=2^n$. Her $n\geq0$ için $n<2^n$ olduğundan $|A|<|\mathcal P(A)|$.

\item Tam sayı katsayılı polinomlar, sonlu tam sayı dizileriyle kodlanabildiği için sayılabilirdir. Her sıfırdan farklı polinomun sonlu sayıda kökü vardır; dolayısıyla cebirsel sayıların kümesi sayılabilir bir sonlu kümeler birleşimidir ve sayılabilirdir. $\mathbb R$ sayılamaz olduğundan geriye kalan aşkın sayılar sayılamazdır.

\item $f:[0,1]\to[0,2]$, $f(x)=2x$ fonksiyonu birebir ve örtendir; tersi $f^{-1}(y)=y/2$'dir.

\item Oteldeki eski $n$ numaralı konuk $2n$ numaralı odaya taşınsın. Böylece bütün tek numaralı odalar boşalır; yeni $k$ numaralı konuk $2k-1$ numaralı odaya yerleştirilir. Hem eski hem yeni bütün konuklar ayrı bir odaya yerleşir.
\end{enumerate}
